% ----------------------------------------------------------------------
% Chapter 4: Approximating hedging strategies by neural networks.
% Assembled from the artifacts
%   lemma4_1_monotone.tex   (monotone sandwich, end of Section 4.1)
%   prop4_3_finite.tex      (finite-space proposition + what breaks, Section 4.2)
%   prop4_7_counterexample.tex (failure on general Omega, Section 4.2)
%   lemma4_4_usc.tex        (usc of OCEs, Section 4.3)
%   lemma4_6_truncation.tex (truncation, Section 4.3)
%   thm4_8_approximation.tex (main theorem, Section 4.3)
% following ch4_structure.md. Old draft labels are preserved in
% "% draft: ..." comments next to each result.
% Remaining external references (Chapters 2 and 3, (A1) to (A6)) are
% hardcoded in the text and marked with "% FAKE REF:" comments; replace
% them by real \ref's at the final merge.
% ----------------------------------------------------------------------
\documentclass[11pt,a4paper]{report}
\usepackage[T1]{fontenc}
\usepackage[utf8]{inputenc}
\usepackage{amsmath,amssymb,amsthm}
\usepackage{enumitem}

% Equation numbers: plain arabic integers, reset each chapter.
% (styling.md item 1 asks for consecutive numbering within each section;
%  with few numbered displays we use chapter-local (1), (2), ... .)
\numberwithin{equation}{chapter}
\renewcommand{\theequation}{\arabic{equation}}

\theoremstyle{plain}
\newtheorem{proposition}{Proposition}[chapter]
\newtheorem{lemma}[proposition]{Lemma}
\newtheorem{theorem}[proposition]{Theorem}
\newtheorem{corollary}[proposition]{Corollary}
\theoremstyle{definition}
\newtheorem{assumption}[proposition]{Assumption}
\newtheorem{definition}[proposition]{Definition}
\newtheorem{example}[proposition]{Example}
% Upright body, bold head (not the italic amsthm "remark" style).
\newtheorem{remark}[proposition]{Remark}

\newcommand{\E}{\mathbb{E}}
\newcommand{\PP}{\mathbb{P}}
\newcommand{\R}{\mathbb{R}}
\newcommand{\N}{\mathbb{N}}
\newcommand{\F}{\mathcal{F}}
\newcommand{\HH}{\mathcal{H}}
\newcommand{\Hu}{\mathcal{H}^{\mathrm{u}}}
\newcommand{\NNet}{\mathcal{N\!N}}
\newcommand{\as}{\ \PP\text{-a.s.}}

\begin{document}

\setcounter{chapter}{3}

\chapter{Approximating hedging strategies by neural networks}
\label{ch:approximation}

% FAKE REF: "(A1)" to "(A6)" throughout this chapter -> Chapter 3,
%   Section 3.4 (full statements of the standing hypotheses).
The previous chapter established that, under the hypotheses (A1) to
(A6), the hedging functional $\pi$ admits an optimal strategy on a
general probability space. This chapter turns to the second question:
can the optimal hedge be approximated on a general $\Omega$ by neural
network strategies? On a finite probability space, B\"uhler et al.\
\cite[Proposition~4.3]{buehler} prove that the network values
$\pi_M(X)$ decrease to $\pi(X)$ as the network capacity $M$ grows.
We identify where that argument uses finiteness of $\Omega$, show by
counterexample that the result fails in general.
We then restore
$\pi_M\to\pi$ under (A1) and (A6), the essential hypotheses of Chapter~3 to prove attainment. Combined with the attainment
theorem of Chapter~3, this yields $\varepsilon$-optimal network
strategies for an attained optimum.

\section{Neural networks, market setup and the approximate problem}
\label{sec:setup}

\subsection*{The market and the hedging functional}

% FAKE REF: market objects (I, S, Hu, H o delta, C_T, W, rho_ell)
%   -> Chapter 2; pi -> Chapter 3; (A1)--(A6) -> Chapter 3, Section 3.4.
We continue in the market setting of Chapter~2 (information process
$I$, mid-price $S$, unconstrained strategies $\Hu$, constraints via
$H\circ\delta$, costs $C_T$ and gains $W$) and write
\begin{equation}\label{eq:pi}
  \pi(X):=\inf_{\delta\in\Hu}\rho\bigl(X+W(H\circ\delta)\bigr)
\end{equation}
for the hedging functional of Chapter~3. Until the repair we do not
impose (A1) to (A6); those hypotheses enter only when we restore
$\pi_M\to\pi$ on general $\Omega$.

\subsection*{Neural networks}

A feed forward neural network with activation function
$\sigma\colon\R\to\R$ is a map
$F=W_L\circ F_{L-1}\circ\dots\circ F_1\colon\R^{d_0}\to\R^{d_1}$ with
$F_j=\sigma\circ W_j$ for affine maps $W_j$, where $\sigma$ acts
componentwise; see B\"uhler et al.\ \cite[Definition~2]{buehler}. We
denote by $\NNet_{\infty,d_0,d_1}$ the set of all such networks and
fix, as in \cite[Section~4.1]{buehler}, an increasing exhaustion
$\NNet_{M,d_0,d_1}\subseteq\NNet_{M+1,d_0,d_1}$ with
$\bigcup_{M\in\N}\NNet_{M,d_0,d_1}=\NNet_{\infty,d_0,d_1}$, each
$\NNet_{M,d_0,d_1}$ parametrized by finitely many weights and
containing the zero network, and such that networks with fewer inputs
embed by zero weight connections. Throughout, the activation function
$\sigma$ is measurable, bounded and nonconstant. The analytical basis
of the approximation theory is the following universal approximation
theorem of Hornik \cite[Theorem~1]{hornik}.

% draft: Hornik's theorem (stated for later reference as thm:hornik).
\begin{theorem}\label{thm:hornik}
  Let $\sigma\colon\R\to\R$ be measurable, bounded and nonconstant, and
  let $d_0\in\N$. Then for every finite measure $\mu$ on $\R^{d_0}$ and
  every $p\in[1,\infty)$, the set $\NNet_{\infty,d_0,1}$ is dense in
  $L^p(\R^{d_0},\mu)$. In particular, applying the result
  componentwise, $\NNet_{\infty,d_0,d}$ is dense in
  $L^p(\R^{d_0},\mu;\R^d)$ for every $d\in\N$.
\end{theorem}

\noindent
The deep hedging approach replaces the infinite-dimensional strategy
class $\Hu$ by finite-dimensional subclasses of network strategies. For
each capacity $M\in\N$, following B\"uhler et al.\
\cite[proof of Proposition~4.3]{buehler}, we set
\[
  \HH_M:=\bigl\{\delta\in\Hu:\delta_k=F_k(I_0,\dots,I_k)
  \text{ with }F_k\in\NNet_{M,r(k+1),d}\bigr\}.
\]
Thus a strategy belongs to $\HH_M$ if and only if each hedge ratio
$\delta_k$ is given by a network of complexity at most $M$ applied to
the information $(I_0,\dots,I_k)$ available at time $t_k$. We do not
include the previous position $\delta_{k-1}$ as an additional network
input. Indeed, each $\delta_k$ is $\F_k$-measurable with
$\F_k=\sigma(I_0,\dots,I_k)$, and $\delta_{k-1}$ is a function of
$(I_0,\dots,I_{k-1},\delta_{k-2})$; recursively, every $\delta_k$ is
already a function of $(I_0,\dots,I_k)$ alone, so feeding
$\delta_{k-1}$ back would not enlarge the class of adapted strategies.
The associated approximate hedging problem is
\begin{equation}\label{eq:piM}
  \pi_M(X):=\inf_{\delta\in\HH_M}\rho\bigl(X+W(H\circ\delta)\bigr).
\end{equation}
The question of this chapter is whether $\pi_M(X)$ approaches $\pi(X)$
as $M\to\infty$. The following elementary observation reduces that
question to a one-sided estimate.

% draft: Lemma 4.1 (monotone sandwich for pi_M);
% same statement as lem:monotone in extension_of_approximation_on_general.tex
% and in the artifact lemma4_1_monotone.tex.
\begin{lemma}\label{lem:monotone}
  For every $X$ and every $M\in\N$, one has
  \[
    \pi_M(X)\ge\pi_{M+1}(X)\ge\pi(X).
  \]
  In particular, $\lim_{M\to\infty}\pi_M(X)$ exists in
  $[-\infty,\infty]$ and dominates $\pi(X)$.
\end{lemma}

\begin{proof}
  By definition, strategies in $\HH_M$ are processes
  $\delta=(\delta_k)_{k=0,\dots,n-1}$ with
  \[
    \delta_k=F_k(I_0,\dots,I_k),\qquad k=0,\dots,n-1,
  \]
  for network maps $F_k$.  Therefore, each $\delta_k$ is $\F_k$-measurable
  and $\delta$ is adapted. With the nested exhaustion of the network
  classes one has $\HH_M\subseteq\HH_{M+1}\subseteq\Hu$, so the infima
  defining $\pi_M$, $\pi_{M+1}$ and $\pi$ in \eqref{eq:pi} and
  \eqref{eq:piM} are taken over increasing sets, which gives both
  inequalities. The limit exists because the sequence
  $(\pi_M(X))_{M\in\N}$ is nonincreasing and bounded below by $\pi(X)$
  in $[-\infty,\infty]$.
\end{proof}

It remains to decide whether the limit equals $\pi(X)$.

\section{The finite-space theorem and what breaks}
\label{sec:finite}

B\"uhler et al., Proposition~4.3, prove the following on a finite
probability space. In the proof we flag each use of finiteness of
$\Omega$ as (F1) to (F3); these are exactly the points that need
attention on a general probability space.

% draft: Proposition 4.3 (Buehler et al.);
% reproduced from the artifact prop4_3_finite.tex.
\begin{proposition}\label{prop:finitespace}
  Suppose the probability space $\Omega=\{\omega_1,\dots,\omega_N\}$ is
  finite with $\PP[\{\omega_i\}]>0$ for all $i$, and $\rho$ is a
  convex risk measure. Then for every claim $X$ one has
  $\displaystyle\lim_{M\to\infty}\pi_M(X)=\pi(X)$.
\end{proposition}

\begin{proof}
  By monotonicity of $(\pi_M)_{M\in\N}$ (Lemma~\ref{lem:monotone}) it
  suffices to show that for every $\varepsilon>0$ some $M$ satisfies
  $\pi_M(X)\le\pi(X)+\varepsilon$.

  Choose $\delta\in\Hu$ with
  \begin{equation}\label{eq:nearopt}
    \rho\bigl(X+W(H\circ\delta)\bigr)\le\pi(X)+\frac{\varepsilon}{2}.
  \end{equation}
  Write $\delta_k=f_k(I_0,\dots,I_k)$ for Borel maps
  $f_k:\R^{r(k+1)}\to\R^d$.

  % Run-in italic labels (not \paragraph): ordinary paragraphs => indent, no extra skip.
  % \indent: amsthm proof often suppresses the first-line indent.
  \indent\textit{(F1) Integrability of the hedge ratios.}
  Because $\Omega$ is finite, each $\delta_k$ takes finitely many values
  and every component of $f_k$ lies in $L^1(\mu_k)$ where $\mu_k$ is the
  law of $(I_0,\dots,I_k)$. Finally, Theorem~\ref{thm:hornik} yields
  networks $F_{k,m}\to f_k$ in $L^1(\mu_k)$.

  \indent\textit{(F2) From $L^1$ to pointwise convergence.}
  On a finite $\Omega$, $L^1$ convergence implies pointwise convergence:
  $\delta^m_k:=F_{k,m}(I_0,\dots,I_k)\to\delta_k$ for every
  $\omega\in\Omega$.

  \indent Set $G_m:=((H\circ\delta^m)\cdot S)_T$ and
  $C_m:=C_T(H\circ\delta^m)$. Since $H$ and the stochastic integral are
  continuous, $G_m\to G$ pointwise. Furthermore, using upper
  semicontinuity of the $c_k$ and continuity of $H$ again, one has
  $\limsup_m C_m=\overline C\le C:=C_T(H\circ\delta)$ pointwise.

  \indent\textit{(F3) Continuity of the risk measure.}
  On the finite space $\Omega$, every random variable may be identified
  with a vector in $\R^N$, so $\rho$ may be viewed as a convex map
  $\R^N\to\R$ and is therefore continuous. Continuity and
  monotonicity of $\rho$ give
  \[
    \lim_{m\to\infty}\rho(X+G_m-C_m)=\rho(X+G-\overline C)
    \le\rho(X+G-C)\le\pi(X)+\frac{\varepsilon}{2}
  \]
  by~\eqref{eq:nearopt}. Pick $m$ large enough so that the
  left-hand side is $\le\pi(X)+\varepsilon$. Since the network families
  are nested with
  $\bigcup_{M\in\N}\NNet_{M,r(k+1),d}=\NNet_{\infty,r(k+1),d}$, there
  exists $M_m$ such that $\delta^m\in\HH_M$ for all $M\ge M_m$, and
  therefore $\pi_M(X)\le\pi(X)+\varepsilon$ for all such $M$.
\end{proof}

\subsection*{What breaks on a general probability space}

\begin{remark}\label{rem:F1}
  \textbf{(F1): hedge ratios need not be integrable.}
  On general $\Omega$, the representing functions $f_k$ need not lie in
  $L^1(\mu_k)$, and then Theorem~\ref{thm:hornik} cannot even be
  invoked: $f_k$ is not an element of the space in which networks are
  dense.
\end{remark}

\begin{example}\label{ex:integrability}
  Let $I_0$ have a standard Cauchy law $\mu_0$ and $\F_0=\sigma(I_0)$.
  The strategy $\delta_0=f(I_0)$ with $f(x)=|x|$ is perfectly admissible
  as a measurable adapted process, but $f\notin L^1(\mu_0)$, so no
  sequence of $\mu_0$-integrable functions (in particular no sequence
  of networks with bounded activation, which are bounded functions)
  converges to $f$ in $L^1(\mu_0)$.
\end{example}

\begin{remark}\label{rem:F2}
  \textbf{(F2): only subsequential almost sure convergence (harmless).}
  On general $\Omega$, $L^1(\PP)$-convergence yields almost sure
  convergence only along a subsequence. This is harmless for the
  argument: the proof only requires \emph{one} good approximating
  strategy, so we may freely pass to subsequences.
\end{remark}

\begin{remark}\label{rem:F3}
  \textbf{(F3): the risk measure is not continuous along almost sure
    convergence.}
  This is the decisive failure. On $\R^N$, convexity of $\rho$ implies
  continuity. On an infinite-dimensional domain, a convex risk measure
  need not be continuous along almost sure convergence. The direction
  needed in the final step of the proof is an \emph{upper}
  semicontinuity:
  \[
    \limsup_{m\to\infty}\rho(Y_m)\le\rho(Y)\qquad\text{whenever }Y_m\to Y\as
  \]
  % FAKE REF: "(A3)" -> the half-bounded Fatou property of Chapter 3,
  %   Section 3.4; lower semicontinuity, i.e. the opposite direction.
  This fails even for the entropic risk measure, which satisfies the
  half-bounded \emph{lower} semicontinuity (A3) of Chapter~3.
\end{remark}

\begin{example}\label{ex:explosion}
  Work on $\Omega=(0,1)$ with the Borel $\sigma$-field and Lebesgue
  measure, let $\lambda>0$ and let
  $\rho_\lambda(Y)=\frac{1}{\lambda}\log\E[e^{-\lambda Y}]$ be the
  entropic risk measure. Set $A_m:=(0,e^{-m})$ and
  \[
    Y_m:=-\frac{2m}{\lambda}\mathbf{1}_{A_m}.
  \]
  For every $\omega\in(0,1)$ one has $\omega\notin A_m$ for all
  $m>\log(1/\omega)$, so $Y_m\to0$ everywhere. But
  \[
    \E\bigl[e^{-\lambda Y_m}\bigr]=1-e^{-m}+e^{-m}e^{2m}
    =1-e^{-m}+e^{m}\longrightarrow\infty,
  \]
  so $\rho_\lambda(Y_m)\to\infty$ while $\rho_\lambda(0)=0$. Small
  events carrying large losses make the risk explode along an almost
  surely vanishing sequence.
\end{example}

\subsection*{Failure on a general probability space}

We now show that Proposition~\ref{prop:finitespace} does not extend to a general probability space. The example is a one-period
market with a perfect hedge whose hedge ratio is
integrable but has no exponential moments.

% draft: Proposition 4.7 (failure without extra assumptions);
% same construction as prop:counterexample in
% extension_of_approximation_on_general.tex and in the artifact
% prop4_7_counterexample.tex.
\begin{proposition}\label{prop:counterexample}
  There exist a one-period market on a general probability space,
  without transaction costs and without trading constraints, an entropic
  risk measure $\rho_\lambda$ with arbitrary risk aversion $\lambda>0$
  and a position $X=-Z$ with $Z\in L^1$ such that
  \[
    \pi(X)\le0\qquad\text{and}\qquad\pi_M(X)=+\infty\quad\text{for every }M\in\N.
  \]
  Moreover, $\rho_\lambda(X+\delta_0(S_1-S_0))=+\infty$ for \emph{every}
  strategy of the form $\delta_0=F(I_0)$ with $F:\R\to\R$ bounded on
  $(0,1)$; this covers all networks with bounded activation and all
  networks with continuous activation.
\end{proposition}

\begin{proof}
  Let $\Omega:=(0,1)\times\{-1,+1\}$ with
  $\PP:=\mathrm{Leb}\otimes(\frac{1}{2}\delta_{-1}+\frac{1}{2}\delta_{+1})$,
  and write $U(\omega):=\omega_1$ and $\xi(\omega):=\omega_2$, so that
  $U$ is uniform on $(0,1)$, $\PP[\xi=\pm1]=\frac{1}{2}$ and $U,\xi$ are
  independent. Take one trading period ($n=1$), information $I_0:=U$,
  $I_1:=\xi$, hence $\F_0=\sigma(U)$ and $\F_1=\sigma(U,\xi)$, one asset
  ($d=1$) with
  \[
    S_0:=1,\qquad S_1:=1+\frac{\xi}{2},
  \]
  no costs ($c_k\equiv0$) and no constraints ($H_0(\delta_0):=\delta_0$).
  Note that $S$ is a $\PP$-martingale and $S_1>0$. Define
  \[
    h(u):=\frac{2}{\lambda}\log\frac{1}{u},\qquad u\in(0,1),\qquad
    Z:=h(U)\frac{\xi}{2},\qquad X:=-Z.
  \]
  Then $\E[|Z|]=\frac{1}{2}\E[h(U)]=\frac{1}{\lambda}<\infty$, so
  $Z\in L^1$.

  \textit{The perfect hedge.}
  The strategy $\delta_0:=h(U)$ is $\F_0$-measurable, and
  \[
    X+\delta_0(S_1-S_0)=-h(U)\frac{\xi}{2}+h(U)\frac{\xi}{2}=0,
  \]
  so $\pi(X)\le\rho_\lambda(0)=0$. Note that $\delta_0=h(U)\in L^1$; the
  optimal strategy is integrable, so Theorem~\ref{thm:hornik} applies to
  it and networks approximate it in $L^1(\PP)$ and, along a subsequence,
  almost surely.

  \textit{Every locally bounded strategy has infinite risk.}
  Let $\delta_0=F(U)$ with $m:=\sup_{u\in(0,1)}|F(u)|<\infty$. Then
  \[
    \rho_\lambda\bigl(X+F(U)(S_1-S_0)\bigr)
    =\frac{1}{\lambda}\log\E\Bigl[\exp\Bigl(\frac{\lambda}{2}\bigl(h(U)-F(U)\bigr)\xi\Bigr)\Bigr].
  \]
  The integrand is nonnegative, so the expectation is at least as large
  as its restriction to the event $\{\xi=+1\}\cap\{h(U)>m\}$. On that
  event one has $h(U)-F(U)\ge h(U)-m>0$, hence
  \begin{align*}
    \exp\Bigl(\frac{\lambda}{2}\bigl(h(U)-F(U)\bigr)\xi\Bigr)
     & =\exp\Bigl(\frac{\lambda}{2}\bigl(h(U)-F(U)\bigr)\Bigr) \\
     & \ge\exp\Bigl(\frac{\lambda}{2}\bigl(h(U)-m\bigr)\Bigr)
    =\frac{e^{-\lambda m/2}}{U}.
  \end{align*}
  Independence of $U$ and $\xi$ and $\PP[\xi=+1]=\frac{1}{2}$ therefore
  yield, with $u_m:=e^{-\lambda m/2}$,
  \[
    \E\Bigl[\exp\Bigl(\frac{\lambda}{2}\bigl(h(U)-F(U)\bigr)\xi\Bigr)\Bigr]
    \ge\frac{e^{-\lambda m/2}}{2}\int_0^{u_m}\frac{1}{u}\,du=\infty,
  \]
  so the entropic risk is $+\infty$.

  \textit{Networks are locally bounded.}
  If the activation $\sigma$ is bounded, every network
  $F=W_L\circ F_{L-1}\circ\dots\circ F_1$ is bounded on all of $\R$,
  because the output of the last hidden layer lies in a bounded set and
  $W_L$ is affine. If $\sigma$ is continuous, $F$ is continuous on $\R$
  and therefore bounded on the compact set $[0,1]\supseteq(0,1)$. In
  either case, every network strategy $\delta_0=F(I_0)$ has
  $\sup_{(0,1)}|F|<\infty$, so $\rho_\lambda(X+W(\delta))=+\infty$ and
  $\pi_M(X)=+\infty$ for every $M$.
\end{proof}

\section{Repair under the Chapter 3 hypotheses}
\label{sec:repair}

The finite-space approximation fails on general $\Omega$ at (F1) and
(F3). We restore it under two hypotheses from Chapter~3: (A1), that the
liability is essentially bounded ($Z\in L^\infty$), and (A6), that every
admissible gain is nonnegative ($W(\delta)\ge0$ a.s.\ for all
$\delta\in\HH$).

% draft: Lemma 4.4 (constant-envelope upper semicontinuity of OCEs);
% same statement as lem:usc-A in extension_with_A1_to_A6.tex and in the
% artifact lemma4_4_usc.tex.
\begin{lemma}\label{lem:usc}
  Let $\ell:\R\to\R$ be convex and nondecreasing and let $\rho_\ell$ be
  the associated OCE. Let $(Y_m)_{m\in\N}$ and $Y$ be random variables
  with
  \[
    Y_m\ge-B\as\text{ for all }m,
  \]
  for some deterministic $B\ge0$. If $\liminf_{m\to\infty}Y_m\ge Y$
  $\PP$-a.s., then
  \[
    \limsup_{m\to\infty}\rho_\ell(Y_m)\le\rho_\ell(Y).
  \]
\end{lemma}

\begin{proof}
  Fix $w\in\R$. Since $-(Y_m+w)\le B+|w|$ and $\ell$ is nondecreasing,
  \[
    g_w:=\ell(B+|w|)\ge\ell\bigl(-(Y_m+w)\bigr),
  \]
  so $g_w$ is a deterministic (hence integrable) constant dominating
  $\ell\bigl(-(Y_m+w)\bigr)$. Extend $\ell$ to
  $[-\infty,\infty)$ by $\ell(-\infty):=\lim_{t\to-\infty}\ell(t)$; the
  extension is continuous and nondecreasing. For any real sequence
  $(a_m)_{m\in\N}$ with $\sup_m a_m<\infty$ one has
  $\limsup_m\ell(a_m)\le\ell(\limsup_m a_m)$. Applying this with
  $a_m=-(Y_m+w)$ and using
  $\limsup_m\bigl(-(Y_m+w)\bigr)=-\liminf_m(Y_m+w)\le-(Y+w)$ a.s.\ yields
  \[
    \limsup_{m\to\infty}\ell\bigl(-(Y_m+w)\bigr)\le\ell\bigl(-(Y+w)\bigr)\as.
  \]
  Reverse Fatou applied to the sequence
  \[
    g_w\ge\ell\bigl(-(Y_m+w)\bigr)
  \]
  gives
  \[
    \limsup_{m\to\infty}\E\bigl[\ell\bigl(-(Y_m+w)\bigr)\bigr]
    \le\E\bigl[\ell\bigl(-(Y+w)\bigr)\bigr].
  \]
  Since $\rho_\ell(Y_m)\le w+\E[\ell(-(Y_m+w))]$, one obtains
  $\limsup_m\rho_\ell(Y_m)\le w+\E[\ell(-(Y+w))]$ for every $w$, and
  taking the infimum over $w$ completes the proof.
\end{proof}

The failure point (F1) is the integrability of the hedge ratios: on a
general probability space, near-optimal strategies need not have
integrable hedge ratios, so Theorem~\ref{thm:hornik} cannot be applied
directly (Example~\ref{ex:integrability}). Under (A1) and (A6), every
admissible objective satisfies
$X+W(\eta)=-Z+W(\eta)\ge-\|Z\|_\infty$, which is the deterministic
lower bound needed for Lemma~\ref{lem:usc}. The next lemma exploits
this to show that bounded strategies are value-dense: every strategy
can be replaced by a bounded one without worsening the OCE objective by
more than an arbitrary $\varepsilon>0$. In particular, $\pi$ equals the
infimum of the objective over bounded strategies alone.

% draft: Lemma 4.6 (truncation / value-density of bounded strategies);
% same statement as lem:truncation-A in extension_with_A1_to_A6.tex and
% in the artifact lemma4_6_truncation.tex.
\begin{lemma}\label{lem:truncation}
  Assume (A1) and (A6), and let $X=-Z$. Then for every $\delta\in\Hu$ and
  every $\varepsilon>0$ there exists a bounded $\delta'\in\Hu$ with
  \[
    \rho_\ell\bigl(X+W(H\circ\delta')\bigr)
    \le\rho_\ell\bigl(X+W(H\circ\delta)\bigr)+\varepsilon.
  \]
  In particular,
  \[
    \pi(X)=\inf\bigl\{\rho_\ell\bigl(X+W(H\circ\delta)\bigr):
    \delta\in\Hu,\ \delta\text{ bounded}\bigr\}.
  \]
\end{lemma}

\begin{proof}
  Fix $\delta\in\Hu$ and write $\eta:=H\circ\delta$, $Y:=X+W(\eta)$.
  Define the componentwise truncations
  $\delta^{(j)}=(\delta^{(j)}_k)_{k=0,\dots,n-1}$ for $j\in\N$ by
  \[
    \delta^{(j)}_k:=
    \begin{cases}
      \delta_k & \text{if }\delta_k\in[-j,j], \\
      -j       & \text{if }\delta_k<-j,       \\
      j        & \text{if }\delta_k>j,
    \end{cases}
    \qquad k=0,\dots,n-1
  \]
  (applied componentwise if $d>1$). Then each $\delta^{(j)}$ is adapted
  (hence in $\Hu$) and bounded by $j$. Set
  $\eta^{(j)}:=H\circ\delta^{(j)}$ and $Y_j:=X+W(\eta^{(j)})$.

  Since there are finitely many time steps, there exists a (random)
  index $j_0<\infty$, namely
  \[
    j_0
    :=
    \max_{\substack{0\le k\le n-1\\ 1\le i\le d}}
    \bigl\lceil\bigl|\delta_k^i\bigr|\bigr\rceil,
  \]
  such that $\delta^{(j)}_k=\delta_k$ for all $k$ and all $j\ge j_0$. In
  particular $\delta^{(j)}_k\to\delta_k$ pointwise for every $k$. The
  recursive definition of $H\circ{}$ then yields $\eta^{(j)}\to\eta$
  pointwise, and therefore $Y_j\to Y$ pointwise.

  By (A6), $W(\eta^{(j)})\ge0$ a.s., so $Y_j\ge-\|Z\|_\infty$ a.s.\ by
  (A1). Lemma~\ref{lem:usc} with $B=\|Z\|_\infty$ gives
  \[
    \limsup_{j\to\infty}\rho_\ell(Y_j)\le\rho_\ell(Y).
  \]
  Choose $j$ large enough that
  $\rho_\ell(Y_j)\le\rho_\ell(Y)+\varepsilon$, and set
  $\delta':=\delta^{(j)}$.
\end{proof}

The following theorem combines the two repairs with the monotone
sandwich of Lemma~\ref{lem:monotone}. It is the main result of this
chapter.

% draft: Theorem 4.8 (approximation under (A1) and (A6));
% same statement as thm:main-A in extension_with_A1_to_A6.tex and in the
% artifact thm4_8_approximation.tex.
\begin{theorem}\label{thm:approximation}
  Assume the standing market and network setting of B\"uhler et al., let
  $\rho=\rho_\ell$ be an OCE risk measure, and assume (A1) and (A6). Then
  \[
    \lim_{M\to\infty}\pi_M(-Z)=\pi(-Z)
  \]
  in $[-\infty,\infty]$, the convergence being monotone from above.
\end{theorem}

\begin{proof}
  Write $X:=-Z$. We distinguish the three cases
  $\pi(X)=+\infty$, $\pi(X)\in\R$ and $\pi(X)=-\infty$.

  \textit{Case $\pi(X)=+\infty$.}
  By Lemma~\ref{lem:monotone}, $\pi_M(X)\ge\pi(X)=+\infty$ for every
  $M$, so $\pi_M(X)=+\infty$ for all $M$ and the claim is trivial.

  \textit{Case $\pi(X)\in\R$.}
  By Lemma~\ref{lem:monotone} it is enough to show that for every
  $\varepsilon>0$ there exists $M$ with
  $\pi_M(X)\le\pi(X)+\varepsilon$.

  \textit{Step 1: a bounded near-optimal target.}
  By Lemma~\ref{lem:truncation} there exists a bounded strategy
  $\delta\in\Hu$, say $|\delta_k|\le j$ a.s.\ for all $k$, such that
  \[
    \rho_\ell\bigl(X+W(H\circ\delta)\bigr)\le\pi(X)+\frac{\varepsilon}{2}.
  \]
  Write $\eta:=H\circ\delta$ and $Y:=X+W(\eta)$.

  \textit{Step 2: representing maps and integrability.}
  Each $\delta_k$ is $\F_k$-measurable with
  $\F_k=\sigma(I_0,\dots,I_k)$, so there exist Borel maps
  $f_k:\R^{r(k+1)}\to\R^d$ with $\delta_k=f_k(I_0,\dots,I_k)$.
  Truncating each component of $f_k$ to the interval $[-j,j]$ changes
  nothing $\PP$-a.s.\ and makes $f_k$ bounded everywhere. In particular
  every component of $f_k$ lies in $L^1(\mu_k)$, where $\mu_k$ is the
  law of $(I_0,\dots,I_k)$. This settles (F1).

  \textit{Step 3: network approximation.}
  By Theorem~\ref{thm:hornik} applied componentwise, there exist
  networks $F_{k,m}\in\NNet_{\infty,r(k+1),d}$ with $F_{k,m}\to f_k$
  in $L^1(\mu_k)$ as $m\to\infty$. Set
  \[
    \delta^m:=\bigl(F_{k,m}(I_0,\dots,I_k)\bigr)_{k=0,\dots,n-1}.
  \]
  Then $\delta^m_k\to\delta_k$ in $L^1(\PP)$ for every $k$. Passing to a
  subsequence (not relabelled) we may assume
  \[
    \delta^m_k\longrightarrow\delta_k\as\qquad\text{for all }k=0,\dots,n-1
  \]
  simultaneously. Since each $F_{k,m}$ belongs to
  $\NNet_{\infty,r(k+1),d}$ and these families are nested with
  $\bigcup_M\NNet_{M,r(k+1),d}=\NNet_{\infty,r(k+1),d}$, there exist
  $M_m\uparrow\infty$ such that $\delta^m\in\HH_{M_m}$.

  \textit{Step 4: constrained strategies and objectives.}
  Set $\eta^m:=H\circ\delta^m$. Continuity of each $H_k(\omega,\cdot)$
  and induction over $k$ give $\eta^m_k\to\eta_k$ $\PP$-a.s.\ for every
  $k$. Gains converge: $(\eta^m\cdot S)_T\to(\eta\cdot S)_T$ a.s.\ Upper
  semicontinuity of the $c_k(\omega,\cdot)$ yields
  \[
    \limsup_{m\to\infty}C_T(\eta^m)\le C_T(\eta)\as,
  \]
  and therefore, writing $Y_m:=X+W(\eta^m)$,
  \[
    \liminf_{m\to\infty}Y_m\ge Y\as.
  \]
  Furthermore, by (A6) and (A1), $Y_m\ge-\|Z\|_\infty$ a.s.\ for every
  $m$.

  \textit{Step 5: conclusion.}
  All conditions of Lemma~\ref{lem:usc} are now satisfied with
  $B=\|Z\|_\infty$, and it gives
  $\limsup_m\rho_\ell(Y_m)\le\rho_\ell(Y)\le\pi(X)+\varepsilon/2$.
  Finally, choose $m_0$ with $\rho_\ell(Y_{m_0})\le\pi(X)+\varepsilon$.
  Since $\delta^{m_0}\in\HH_{M_{m_0}}$, one has
  $\pi_M(X)\le\pi(X)+\varepsilon$ for all $M\ge M_{m_0}$.

  \textit{Case $\pi(X)=-\infty$.}
  Let $L\in\R$ be arbitrary. Choose $\delta\in\Hu$ with
  $\rho_\ell(X+W(H\circ\delta))\le L-1$. Then, steps 1 to 5 above go
  through with the threshold $L$ in place of $\pi(X)+\varepsilon$,
  yielding $M$ with $\pi_M(X)\le L$. By choosing $L$ lower and lower,
  and being $(\pi_M(X))_{M\in\N}$ nonincreasing,
  $\pi_M(X)\to-\infty=\pi(X)$.
\end{proof}

% draft: Remark after Theorem 4.8 (what is not used; weakening (A1));
% condensed from rem:unused and rem:A1prime in extension_with_A1_to_A6.tex.
\begin{remark}\label{rem:unused}
  The proof uses (A1) and (A6), the standing continuity structure of
  the maps $H_k$ and cost functions $c_k$, the OCE representation of
  $\rho$ and Theorem~\ref{thm:hornik}. It does not use (A2) to (A5)
  (convexity of $\HH$, the half-bounded Fatou property,
  $L^0$-boundedness of the gain set and closedness of the dominated
  gain cone). Those hypotheses
  remain needed for attainment in Chapter~3; they are not needed for
  value convergence. The full strength of (A1) is also not needed: the
  same proof works if $Z$ is merely bounded from below with
  $\E[\ell(Z+c)]<\infty$ for every $c\ge0$, the deterministic envelope
  in Lemma~\ref{lem:usc} being replaced by a random $\ell$-integrable
  one. We do not pursue this weakening.
\end{remark}

Combining Theorem~\ref{thm:approximation} with the attainment theorem
of Chapter~3 gives the punchline of this thesis.

% draft: Corollary 4.10 (approximation of an attained optimum);
% same statement as cor:attained-A in extension_with_A1_to_A6.tex.
\begin{corollary}\label{cor:attained}
  Assume in addition \textup{(A2)} to \textup{(A5)}, so that the
  attainment theorem of Chapter~3 applies. Then there exists
  $\delta^*\in\HH$ with $\rho_\ell(-Z+W(\delta^*))=\pi(-Z)$, and for
  every $\varepsilon>0$ there exist $M\in\N$ and
  $\delta^\varepsilon\in\HH_M$ with
  \[
    \rho_\ell\bigl(-Z+W(H\circ\delta^\varepsilon)\bigr)
    \le\rho_\ell\bigl(-Z+W(\delta^*)\bigr)+\varepsilon.
  \]
\end{corollary}

\begin{proof}
  The existence of $\delta^*$ follows from the attainment theorem of
  Chapter~3, and the $\varepsilon$-statement is
  Theorem~\ref{thm:approximation}.
\end{proof}

In one sentence: under (A1) to (A6), neural network strategies are
$\varepsilon$-optimal for an \emph{attained} optimum.

\section{From theory to algorithm}
\label{sec:algorithm}

Theorem~\ref{thm:approximation} justifies replacing the strategy class
$\Hu$ by the network classes $\HH_M$; this section records why that
replacement makes the hedging problem computable. For an OCE one has
$\rho_\ell(Y)=\inf_{w\in\R}\{w+\E[\ell(-Y-w)]\}$, so with
$Y=-Z+W(H\circ\delta^\theta)$, where $\delta^\theta\in\HH_M$ is
determined by the finitely many network weights $\theta$ of the maps
$F_{\theta_k}$, the approximate problem \eqref{eq:piM} becomes
\[
  \pi_M(-Z)
  =\inf_{(w,\theta)}
  \Bigl\{w+\E\bigl[\ell\bigl(Z-W(H\circ\delta^\theta)-w\bigr)\bigr]\Bigr\}.
\]
This is a finite-dimensional minimization over the pair $(w,\theta)$,
jointly and without any inner optimization. In practice the expectation
is estimated by minibatches of simulated (or historically calibrated)
sample paths of the information process $I$, the gradient in
$(w,\theta)$ is computed by backpropagation through the semi-recurrent
parametrization $\delta_k=F_{\theta_k}(I_0,\dots,I_k,\delta_{k-1})$,
and the minimization is carried out by stochastic gradient descent. For
the entropic risk measure the inner minimizer is explicit,
\[
  w^*=\frac{1}{\lambda}\log\E\bigl[e^{\lambda(Z-W(H\circ\delta^\theta))}\bigr],
\]
so one minimizes over $\theta$ alone; cf.\ B\"uhler et al.\
\cite[Section~3.2]{buehler}.

This thesis contains no numerical experiments of its own. B\"uhler et
al.\ \cite[Section~5]{buehler} train these networks and backtest the
resulting hedges: in their experiments the learned strategies achieve a
substantial reduction of the hedging risk relative to the unhedged
position and remain competitive with classical model-based benchmarks
in the presence of transaction costs, where the classical strategies
are no longer optimal. One structural caveat must be stated explicitly:
the networks are trained exclusively on historic data or on simulated
data calibrated to historic observations. A backtest evaluates the
hedge on the same statistical regime it was trained on, and past
performance is not indicative of future results; this limitation is
inherent to the method and not an implementation detail. Finally, for
general convex risk measures beyond OCEs, the robust representation
leads to a minimax problem over two networks
(B\"uhler et al.\ \cite[Section~4.5]{buehler}); this extension is not
pursued in this thesis.

\section{Strategies need not converge}
\label{sec:strategies}

Theorem~\ref{thm:approximation} and Corollary~\ref{cor:attained} are
statements about \emph{values}: the network values converge to
$\pi(-Z)$, and near-minimizers are $\varepsilon$-optimal. We close the
chapter by observing that nothing of this kind holds for the strategies
themselves.

% draft: Example (nonuniqueness blocks strategy convergence);
% same as ex:nonunique in extension_with_A1_to_A6.tex.
\begin{example}\label{ex:nonunique}
  Let $S\equiv s_0$ be constant, $C_T\equiv0$ and $\HH=\Hu$ (no
  constraints). Then $W(\delta)=0$ for every $\delta$, so every
  strategy is optimal for every $Z$ and $\pi(-Z)=\rho_\ell(-Z)$.
  Network near-minimizers may oscillate indefinitely between distinct
  elements of $\HH_M$ without converging in any topology. The same
  phenomenon occurs with two identical assets and zero costs: the
  strategies $(\delta,0)$ and $(0,\delta)$ produce the same gain, and
  minimizers may alternate between them.
\end{example}

\begin{remark}\label{rem:values-not-strategies}
  Theorem~\ref{thm:approximation} and Corollary~\ref{cor:attained}
  should therefore be read as value statements only. Without strict
  convexity of the loss function $\ell$ and a nonredundancy condition
  on $S$, optimal strategies need not be unique, and a sequence of
  network near-minimizers $\delta^M\in\HH_M$ need not converge to an
  optimal strategy $\delta^*$ in any reasonable topology. What the
  theory guarantees is that the \emph{risk} of $\delta^M$ approaches
  the optimal risk; which minimizer the training procedure selects is
  not determined by the value convergence.
\end{remark}

\begin{thebibliography}{9}

  \bibitem{buehler}
  H. B\"uhler, L. Gonon, J. Teichmann and B. Wood,
  Deep hedging,
  \emph{Quantitative Finance} 19(8), 1271-1291, 2019.

  \bibitem{hornik}
  K. Hornik,
  Approximation capabilities of multilayer feedforward networks,
  \emph{Neural Networks} 4(2), 251-257, 1991.

\end{thebibliography}

\end{document}
